\begin{abstract}
Vision-Language-Action (VLA) models leverage large-scale vision-language pretraining for flexible robot manipulation, yet at test time they remain brittle to changed object positions and to familiar scenes paired with different instructions. 
A growing family of methods addresses this brittleness by supplying the policy with grounding signals, such as 2D pixel coordinates for object localization and placement. 
However, we find that how the grounding signal is represented and injected matters more than the signal itself. In this work, we propose a lightweight module that represents the grounding signal in 3D and injects the resulting embedding directly into the action head. The module is a two-layer MLP and requires no changes to the VLA backbone or pretraining pipeline, yet it yields substantially larger gains than language- or visual-prompting alternatives. 
On LIBERO-PRO, our method improves the average success rate of GR00T-N1.6 from $31.2$ to $77.5$ under task perturbation and from $28.1$ to $60.2$ under position perturbation. Comparable gains are also achieved for $\pi_{0.5}$, demonstrating that the mechanism is backbone-agnostic across VLAs with diffusion-based action heads. 
We further validate the practical applicability with real-world experiments. 
Together, these results support our central finding: lifting adequate 2D grounding into 3D and injecting it into the action head enables spatial and instance-level task generalization in VLAs. 
Project page: \href{https://3dgnd-head-injection.github.io/}{\texttt{3dgnd-head-injection.github.io}}

% Vision-Language-Action (VLA) models leverage large-scale vision-language pretraining for flexible robot manipulation, yet at test time they remain brittle along two axes: spatial generalization, when object positions differ from those seen during training, and task generalization, when a familiar scene is paired with a different language instruction than the one seen in training.
% A growing family of methods addresses this brittleness by endowing a policy with the spatial and task-aware information such as 2D pixel-coordinate for object localization and placement.
% However, we find that existing representation through language prompting or visual prompting does not address the limitations; in contrast, exploiting a 3D point-based representation and feeding it directly to the action head leads to substantial improvements—revealing that how the grounding signal is represented and injected into the VLA is the true game changer.
% Thus, we propose a lightweight, model-agnostic module that represents the grounding signal in 3D, computes its relative displacement to the gripper, and injects the resulting spatial embedding directly into the action head through adaptive layer normalization. The entire module is a two-layer MLP that requires no changes to the VLA backbone or pretraining pipeline.
% On LIBERO-PRO, our method improves the average success rate of GR00T-N1.6 from 31.2 to 77.5 points under task perturbation and from 28.1 to 60.2 points under position perturbation (gains of 46.3 and 32.1 points, respectively). Comparable gains are achieved for $\pi_{0.5}$ as well, demonstrating that the mechanism is backbone-agnostic.
% We further validate practical applicability with real-world experiments. Together, these results support our central finding: given adequate grounding lifted into 3D, injecting it directly into the action head is what unlocks both spatial and task generalization in VLAs—achievable with nothing more than a lightweight module on top of a pretrained backbone. 
\end{abstract}